\documentclass[11pt]{report}
%\usepackage{fourier}
\usepackage[T1]{fontenc}
\usepackage[utf8]{inputenc}

%--------------------------
\usepackage{comment}
\usepackage[top=2cm,bottom=2cm,left=2cm,right=2cm,
                marginparwidth=1.1cm,marginparsep=2mm]{geometry}
%\usepackage{geometry}
\usepackage{tikz}
%\usepackage[]{kpfonts}
\usepackage{amsmath,amssymb,amsfonts,mathtools,mathptmx, amsthm,amsrefs}
\renewcommand{\qedsymbol}{$\mid \mid \mid $}
\usepackage{enumerate}
\usepackage{enumitem,pifont}
%\usepackage{kpfonts}
% Couleurs
\usepackage{mathrsfs}
\usepackage{xcolor}
\definecolor{bleueLeger}{RGB}{70,130,180}
\definecolor{orange}{RGB}{255,138,88}
\definecolor{blue}{RGB}{23,74,117}
\definecolor{vert}{RGB}{21,122,81}
\definecolor{jaune}{RGB}{255,185,88}
\definecolor{grise}{RGB}{200,200,200}

%autres paquets
\usepackage{blindtext}
\usepackage{color}
\usepackage{titlesec}
\usepackage{tocloft}
%Les fontes
\usepackage{anyfontsize}

%les graphiques
\usepackage{graphicx}


%la couleur de fon de la page
\usepackage[angle=45]{background}
\backgroundsetup{contents={voTex}}


% on peut charger ce paquet avec les options suivantes suivantes: pages=some, placement=bottom

%styles des éléments de sectionnement
\usepackage{sectsty}
\sectionfont{\bf\color{black}\sectionrule{3pt}{0pt}{-1ex}{2pt}}
%couleurs
\definecolor{bleue}{RGB}{70,130,180}


%style des chapitres
\renewcommand{\thechapter}{\Roman{chapter}}
\titleformat{\chapter}[display]
            {\bfseries\Large}
            {\filcenter\normalfont\textbf{\chaptertitlename} \large
             \textbf{ \thechapter}}
            {4ex}
            {\color{bleue}%\titlerule[5pt]
              \color{black}\Huge
            % \vspace{2ex}
             \filcenter}
            [\vspace{3ex}
              \color{bleue}\thispagestyle{plain}]
              
%style des sections
     
     
\makeatletter


\usepackage{fancyhdr}
\pagestyle{fancy}
 \fancyfoot{ }
%\footrule{1pt}
%\footrulewidth{1}
\fancyhead[RO]{ }
%\fancyfoot[LO,LE]{\bsc{Nom}}
\fancyfoot[C,C]{\thepage}

\definecolor{beaublue}{rgb}{0.74, 0.83, 0.9}

%\renewcommand{footrulewidth}{1pt} 
%les théorèmes encadrés

\usepackage[most]{tcolorbox}

\newenvironment{boxtheorem}
{\begin{tcolorbox}
[enhanced jigsaw,breakable,pad at break*=2mm,
 colback=black!5,colframe=bleueLeger]}
{\end{tcolorbox}}
\newenvironment{boxdefinition}
{\begin{tcolorbox}
		[enhanced jigsaw,breakable,pad at break*=1mm,
		colback=black!5!white,boxrule=0pt,frame hidden,
		borderline west={1.5mm}{-2mm}{blue}]}
	{\end{tcolorbox}}
\newenvironment{boxcorollaire}
{\begin{tcolorbox}
		[enhanced jigsaw,breakable,pad at break*=1mm,
		colback=black!5!white,boxrule=0pt,frame hidden,
		borderline east={1.5mm}{-2mm}{blue}]}
	{\end{tcolorbox}}


\swapnumbers
\newtheoremstyle{styletheorem}
{0pt}{0pt}{\itshape}{2pt}{\bf\sffamily\color{bleueLeger}}{\;}{0.25em}
{\small\sffamily\color{orange}\thmname{#1}
 \nobreakspace\thmnumber{\@ifnotempty{#1}{}\@upn{#2}}
 \thmnote{\nobreakspace\the\thm@notefont\sffamily\bfseries\color{black} (#3)}}
%%%%styles définitions, remarques,notes, conclusions, etc...
\newtheoremstyle{styledef}
{0pt}{0pt}{\fontshape{it}}{2pt}{\bf\sffamily\color{bleueLeger}}{\;}{0.25em}
{\small\sffamily\color{orange}\thmname{#1}
 \nobreakspace\thmnumber{\@ifnotempty{#1}{}\@upn{#2}}
 \thmnote{\nobreakspace\the\thm@notefont\sffamily\bfseries\color{black} (#3)}}


\theoremstyle{styletheorem}
\newtheorem{theoreme}{Théorème}[section]
%\newtheorem{theorem}{Théorème}[chapter]
\newtheorem{corollary}[theoreme]{Corollaire}
\newtheorem{lemma}[theoreme]{Lemme}
\newtheorem{lm}[theoreme]{Lemme}
%\newtheorem{corollaire}[theoreme]{Corollaire}
\newtheorem{algo}[theoreme]{Algorithme}
\newtheorem{co}[theoreme]{Corollaire}
\newtheorem{conj}[theoreme]{Conjecture}
\newtheorem{proposition}[theoreme]{Proposition}
\newtheorem{remarque}[theoreme]{\textbf{Remarque}}
\newtheorem{rema}[theoreme]{\textbf{Remarques et Notations}}
\newenvironment{theorem}
{\vspace*{0.3cm}\begin{boxtheorem}\begin{theoreme}}{\end{theoreme}\end{boxtheorem}\vspace*{0.3cm}}
\newenvironment{tho}
{\vspace*{0.3cm}\begin{boxtheorem}\begin{theoreme}}{\end{theoreme}\end{boxtheorem}\vspace*{0.3cm}}

%%%%%Styles pour Dzfinition, Condition, Exemples, Algo, Question, Axiommes, Propriété, hypothèses
\theoremstyle{styledef}
\newtheorem{definitions}[theoreme]{Définition}
\newtheorem{exa}[theoreme]{Exemple}
\newtheorem{ex}[theoreme]{Exemples}
\newtheorem{exetco}[theoreme]{Exemple et Contre-exemple}
\newtheorem{propozition}[theoreme]{Proposition-Définition}

\theoremstyle{remark}
%\newtheorem{remarque}[theoreme]{\textbf{Remarque}}
\newtheorem*{nte}{\textbf{Note}}
%\newtheorem{thm}{Th\'{e}or\`{e}me}[chapter] 
%\newtheorem{thms}[subsubsection]{Th\'{e}or\`{e}me}



\newenvironment{lemme}
{\vspace*{0.3cm}\begin{boxtheorem}\begin{lemma}}{\end{lemma}\end{boxtheorem}\vspace*{0.3cm}}
\newenvironment{prop}
{\vspace*{0.3cm}\begin{boxtheorem}\begin{proposition}}{\end{proposition}\end{boxtheorem}\vspace*{0.3cm}}
\newenvironment{ren}
{\vspace*{0.3cm}\begin{boxtheorem}\begin{rema}}{\end{rema}\end{boxtheorem}\vspace*{0.3cm}}
\newenvironment{p}
{\vspace*{0.3cm}\begin{boxtheorem}\begin{proposition}}{\end{proposition}\end{boxtheorem}\vspace*{0.3cm}}
\newenvironment{corollaire}{\begin{boxcorollaire}\begin{corollary}}{\end{corollary}\end{boxcorollaire}}
\newenvironment{definition}{\vspace*{0.3cm}\begin{boxdefinition}\begin{definitions}}{\end{definitions}\end{boxdefinition}\vspace*{0.3cm}}
\newenvironment{rem}{\vspace*{0.3cm}\begin{boxdefinition}\begin{remarque}}{\end{remarque}\end{boxdefinition}\vspace*{0.3cm}}


%%%%%%%%%%%%%%%%%%%%%%%%%%%%%%%%%%%%%%%
%%%%%%%%%%%%%%%%%%%%%%%%%%%%%%%%%%%%%%

\newenvironment{propdf}{\vspace*{0.3cm}\begin{boxdefinition}\begin{propozition}}{\end{propozition}\end{boxdefinition}\vspace*{0.3cm}}







%%%%%%%%%%%%%%%%%%%%%%%%%%%%%%%%%%%%%%%%%
%%%%%%%%%%%%%%%%%%%%%%%%%%%%%%%%%%%%%%%%%%%
%Listes stylées avec le paquet tcolorbox
\newenvironment{monenumerate}{%
\begin{enumerate}}{\end{enumerate}\vspace*{0.3cm}}
\tcolorboxenvironment{monenumerate}{blanker,breakable,pad at break*=1mm,before skip=6pt, after skip=6pt,borderline west={-1mm}{0pt}{bleueLeger}}

\newtheorem{dfs}[theoreme]{D\'{e}finition}
\newtheorem{thme}{Th\'{e}or\`{e}me}[theoreme]


%\numberwithin{equation}{equation}
\renewcommand*{\familydefault}{\sfdefault}
%\newtheorem{prpt}{Part}
%================================================
%===========Set of Functions==================
\newcommand{\app}{\mathrm{App}}
%=============================================


\usepackage[tight,francais]{minitoc}
\usepackage{multicol}
%\usepackage{fancyhdr}
\usepackage{fancybox}
%\renewcommand{\thechapter}{\Roman{chapter}}
%\renewcommand{proof}{\textit{Preuve}}


\usepackage{setspace}
%\renewcommand{\thesection}{\S\Roman{chapter}.\Alph{section}{}{}}
\usepackage{hyperref}
\hypersetup{backref,colorlinks=false,linkcolor=black}

%\hypersetup{backref,pdfpagemode=FullScreen,colorlinks=true,linkcolor=blue}

%%%%%%%%%%%%
%Les scripts
%--------------------------------------------------------------------------

\newcommand{\ie}{c'est-à-dire }
\newcommand{\Ie}{C'est-à-dire }
\newcommand{\ssi}{si et seulement si\  }
\newcommand{\R}{\ensuremath{\mathbb{R}}}
\newcommand{\N}{\ensuremath{\mathbb{N}}}
%\newcommand{\C}{\ensuremath{\mathbb{C}}}
\newcommand{\K}{\ensuremath{\mathbb{K}}}
\newcommand{\ep}{{1\leq i \leq p,1\leq j \leq q}}
\newcommand{\eq}{{1\leq i \leq q,1\leq j \leq p}}
\newcommand{\en}{{1\leq i \leq n,1\leq j \leq m}}
\newcommand{\eij}{{1\leq i,j \leq p}}
\newcommand{\Fe}{\ensuremath{Fe_{\tau}(X)}}
\newcommand{\Fr}{\ensuremath{\mathcal{F}}  }
\newcommand{\Pk}{\ensuremath{P_k}  }
\newcommand{\F}{\ensuremath{\mathcal{F}}}
\newcommand{\f}{\ensuremath{\mathscr{F}}}



%-----------------------------------------------------------------------------------------------------
\usepackage{pifont}
\usepackage[tight]{shorttoc}
\usepackage{lmodern}
\usepackage[french]{babel}
%\usepackage[T1]{fontenc}
\newcommand{\first}[1]{{\color{steelblue}\calligra #1}}
\everymath{\displaystyle}
\usepackage{hyperref}
\usepackage{enumitem}

\usepackage{amsmath, amsfonts, amssymb}





%%%%%%%%%%%%%

%
\newsavebox{\auteurbm}
\newenvironment{Bonmot}[1]%
{\small\slshape%
	\savebox{\auteurbm}{\upshape\sffamily#1}%
	\begin{flushright}}
	{\\[4pt]\usebox{\auteurbm}
	\end{flushright}\normalsize\upshape}
%

\usepackage{mathpazo}
%\renewcommand{\familydefault}{ttdefault}

\newcommand*\closure[1]{\ensuremath{\overline{#1}}}


\newcommand\frontmatter{%
	\cleardoublepage
	%\@mainmatterfalse
	\pagenumbering{roman}}

\newcommand\mainmatter{%
	\cleardoublepage
	% \@mainmattertrue
	\pagenumbering{arabic}}





